% ===================== Chapter 9 =====================
\chapter{Structs and Methods}
\label{ch:structs-methods}

So far our data has been loose: an \code{i32} here, a \code{str} there. But the
things we model are rarely single values. A bank account has an owner \emph{and} a
balance; a point has an $x$ \emph{and} a $y$. A \textbf{struct} lets you bundle
several values into one named type, and \textbf{methods} let you attach behaviour
to that type. Together they are how you build your own data types in Salam.

\section{Defining a struct}

A struct is declared with the \code{struct} keyword, a name, and a list of
\textbf{fields} each a name, a type, and a default value closed with
\code{end}:

\begin{salam}
struct Point:
    x i32 = 0
    y i32 = 0
end
\end{salam}

This defines a new type, \code{Point}, with two fields. Each field has a default
value (here \code{0}), which is used when you do not supply one. \code{Point} is
now a type you can use anywhere a type is expected, just like \code{i32} or
\code{str}.

\section{Creating instances}

To make a value of a struct type an \textbf{instance} write the type name
followed by braces listing the fields you want to set:

\begin{salam}
struct Point:
    x i32 = 0
    y i32 = 0
end

func main:
    origin Point = Point {}                 // both fields default to 0
    p Point = Point { x = 3, y = 4 }        // set both fields
    println origin.x, origin.y
    println p.x, p.y
end
\end{salam}

\begin{progout}
0 0
3 4
\end{progout}

Inside the braces you assign fields by name with \code{field = value}. Any field
you leave out keeps its default \code{Point \{\}} relies on the defaults for
both, giving \code{(0, 0)}. You reach a field of an instance with a dot:
\code{p.x} is \code{p}'s \code{x} field.

\begin{note}
Because every field has a default, \code{Point \{\}} is always a valid, fully
formed value there is no such thing as a half-built struct with missing
fields. This is why field defaults are required in the struct definition: they
guarantee that an instance is never left with an uninitialised field.
\end{note}

\section{Reading and changing fields}

Field access with the dot works both ways: you can read a field, and if the
instance is \code{mut} you can assign to one:

\begin{salam}
struct Point:
    x i32 = 0
    y i32 = 0
end

func main:
    mut p Point = Point { x = 1, y = 1 }
    p.x = 10
    p.y = p.y + 5
    println p.x, p.y
end
\end{salam}

\begin{progout}
10 6
\end{progout}

\section{Structs are copied by value}

A struct behaves like the simple values of Chapter~2: assigning it, passing it to
a function, or returning it makes a \emph{copy}. Changing the copy does not affect
the original:

\begin{salam}
struct Point:
    x i32 = 0
    y i32 = 0
end

func main:
    mut p Point = Point { x = 1, y = 2 }
    mut q Point = p        // q is a copy of p
    q.x = 99
    println p.x, q.x       // p is untouched
end
\end{salam}

\begin{progout}
1 99
\end{progout}

This ``value semantics'' makes structs predictable: a \code{Point} you hand to a
function cannot be quietly modified by that function. It is the same guarantee you
saw for ordinary arguments in Chapter~6, now extended to your own types.

\section{Methods: behaviour on a struct}

A struct can carry not just data but \textbf{methods} functions that belong to
the type and act on a particular instance. You write a method inside the
\code{struct} block, exactly like a function, and refer to the current instance
as \code{this}:

\begin{salam}
struct Counter:
    value i32 = 0
    func increment():
        this.value = this.value + 1
    end
    func get() i32:
        ret this.value
    end
end

func main:
    mut c Counter = Counter {}
    c.increment()
    c.increment()
    c.increment()
    println c.get()
end
\end{salam}

\begin{progout}
3
\end{progout}

You \emph{call} a method through an instance with the dot: \code{c.increment()}.
Inside the method, \code{this} refers to the instance the method was called on ---
so \code{this.value} is \emph{that} counter's \code{value}. Each \code{Counter}
keeps its own \code{value}, and a method always works on the one it was called
through.

\begin{note}
Unlike an ordinary argument, a method \emph{can} change the instance it is called
on: \code{increment} really does update \code{c}'s \code{value}. This is the
sanctioned way to modify data in place that we hinted at in Chapter~6 a
struct's own methods are trusted to change it, while outside functions get
copies.
\end{note}

\section{Methods that compute and methods that act}

Just like free functions, a method may return a value or return nothing. A method
that returns a value is written with a return type after its parameter list:

\begin{salam}
struct Rectangle:
    width i32 = 0
    height i32 = 0
    func area() i32:
        ret this.width * this.height
    end
    func is_square() bool:
        ret this.width == this.height
    end
end

func main:
    r Rectangle = Rectangle { width = 4, height = 4 }
    println "area:", r.area()
    println "square?", r.is_square()
end
\end{salam}

\begin{progout}
area: 16
square? true
\end{progout}

Here \code{area} and \code{is\_square} only \emph{read} the rectangle (through
\code{this}) and compute an answer; they do not change anything. A method may also
take parameters of its own, in addition to the implicit \code{this}:

\begin{salam}
struct Rectangle:
    width i32 = 0
    height i32 = 0
    func scale(factor i32):
        this.width = this.width * factor
        this.height = this.height * factor
    end
end

func main:
    mut r Rectangle = Rectangle { width = 2, height = 3 }
    r.scale(5)
    println r.width, r.height
end
\end{salam}

\begin{progout}
10 15
\end{progout}

\section{A fuller example: a bank account}

Let us combine fields and methods into something realistic a bank account that
guards its own balance. The \code{withdraw} method returns a \code{bool} reporting
whether the withdrawal succeeded:

\begin{salam}
import "str"

struct Account:
    owner str = ""
    balance i32 = 0
    func deposit(amount i32):
        this.balance = this.balance + amount
    end
    func withdraw(amount i32) bool:
        if amount > this.balance:
            ret false
        end
        this.balance = this.balance - amount
        ret true
    end
    func describe() str:
        ret this.owner + " has " + str.FromInt(this.balance as i64)
    end
end

func main:
    mut acc Account = Account { owner = "Sara", balance = 100 }
    acc.deposit(50)
    println acc.describe()
    println "withdraw 30?", acc.withdraw(30)
    println "withdraw 1000?", acc.withdraw(1000)
    println acc.describe()
end
\end{salam}

\begin{progout}
Sara has 150
withdraw 30? true
withdraw 1000? false
Sara has 120
\end{progout}

This is the heart of good design: the \code{Account} \emph{owns} its balance, and
the only way to change it is through \code{deposit} and \code{withdraw}, which
enforce the rules (you cannot withdraw more than you have). The data and the
operations that protect it live together in one place.

\section{Structs within structs}

A field may itself be a struct, letting you build larger structures from smaller
ones. You reach a nested field by chaining dots:

\begin{salam}
struct Point:
    x i32 = 0
    y i32 = 0
end

struct Line:
    start Point = Point {}
    finish Point = Point {}
end

func main:
    ln Line = Line {
        start = Point { x = 0, y = 0 },
        finish = Point { x = 5, y = 5 }
    }
    println ln.start.x, ln.finish.y
end
\end{salam}

\begin{progout}
0 5
\end{progout}

\code{ln.finish.y} reads the \code{y} field of the \code{finish} field of
\code{ln}. Composing structs this way is how you model anything with structure ---
a line from two points, a person with an address, a game piece with a position
and a colour.

\begin{deep}[In Depth: structs become C structs]
Each Salam struct compiles to a plain C \code{struct} with the same fields in the
same order, and a method compiles to an ordinary C function that takes a pointer
to the struct as a hidden first argument which is exactly what \code{this} is.
Because the layout is a direct, flat C struct, accessing a field is a single
memory read with no indirection, and passing a struct by value is a simple
memory copy. This is why structs are both convenient \emph{and} fast, and why
they map so cleanly onto C libraries through the foreign function interface
(Chapter~19).
\end{deep}

\section{Chapter summary}

\begin{itemize}[leftmargin=1.4em]
  \item A \code{struct} bundles named fields, each with a default value, into a
        new type.
  \item Create an instance with \code{Type \{ field = value, ... \}}; omitted
        fields use their defaults.
  \item Read and write fields with the dot: \code{p.x}.
  \item Structs are copied by value on assignment, argument passing, and return.
  \item Methods live inside the \code{struct} block, refer to the instance as
        \code{this}, and are called with the dot: \code{c.increment()}.
  \item A method may read or modify \code{this}, take its own parameters, and
        return a value or nothing.
  \item Struct fields may themselves be structs; chain dots to reach nested
        fields.
\end{itemize}

\section{Exercises}

\exercise Define a \code{struct Circle} with a \code{radius f64} field and a
method \code{area() f64} that returns $\pi r^2$ (use \code{3.14159} for $\pi$).
Print the area of a circle of radius 2.

\exercise Add a \code{describe()} method to \code{Point} that returns a string
like \code{"(3, 4)"}, and print it for a couple of points.

\exercise Build a \code{struct Stack} around a fixed \code{i32[16]} array and a
\code{count}, with \code{push} and \code{pop} methods. Push three values and pop
them back.

\exercise Extend the \code{Account} example with a \code{transfer(other, amount)}
idea: withdraw from one account and deposit into another, only if the withdrawal
succeeds.

\exercise Define a \code{struct Rectangle} with a \code{perimeter()} method as well
as \code{area()}, and print both for a $3 \times 5$ rectangle.

\exercise Model a \code{struct Person} containing a \code{name} and an
\code{Address} struct (street, city). Create one and print the city via chained
field access.
